% https://www.ms.uky.edu/~perry/213-f19-perry/_assets/sources/lec_41_16_9.tex

\documentclass[notes=show,10pt]{beamer}
\usepackage{mathpazo}			%% Beamer stuff
\usepackage{amsmath}			%% AMS math packages
\usepackage{amsfonts}			%%	-	special fonts
\usepackage{amssymb}			%%	-	special symbols
\usepackage{graphicx}			%% graphics
\usepackage{tcolorbox}			%% color boxes

\usepackage{hyperref}			%% for web links
\hypersetup{
%    bookmarks=true,         		% show bookmarks bar?
    unicode=false,          		% non-Latin characters in Acrobat’s bookmarks
    pdftoolbar=true,        		% show Acrobat’s toolbar?
    pdfmenubar=true,        	% show Acrobat’s menu?
    pdffitwindow=false,     		% window fit to page when opened
    pdfstartview={FitH},    		% fits the width of the page to the window
    pdfnewwindow=true,      	% links in new PDF window
    colorlinks=true,       			% false: boxed links; true: colored links
    linkcolor=blue,          		% color of internal links (change box color with linkbordercolor)
    citecolor=PineGreen,    		% color of links to bibliography
    filecolor=magenta,      		% color of file links
    urlcolor=cyan           		% color of external links
}


%%%%%%%%%%%%%%%%%%%%%%%%%%%%%%%%%%%
%
%		TiKZ Ploting Stuff
%
%%%%%%%%%%%%%%%%%%%%%%%%%%%%%%%%%%%

\usepackage{tikz}				%% TikZ -graphics

\usetikzlibrary{arrows, positioning, calc, intersections}
\usetikzlibrary{decorations.pathreplacing, decorations.markings}

\usepackage{pgfplots}			%% fancy plotting

%%%%%%%%%%%%%%%%%%%%%%%%%%%%%%%%%%
%
%	For 3d plotting, using the tikz-3dplot package so
%	that (x,y,z) coordinates are specified in normal order
%	and one can change the perspective for 3-d plots
%
%	To invoke, use
%
%	\begin{tikzpicture}[tdplot_main_coords,{other options}]
%		...
%	\end{tikzpicture}
%
%%%%%%%%%%%%%%%%%%%%%%%%%%%%%%%%%%

\usepackage{tikz-3dplot}		%% Enhanced 3d plotting

%%%%%%%%%%%%%%%%%%%%%%%%%%%%%%%%%%
%
%		Package forest to draw tree diagrams for chain rule
%
%%%%%%%%%%%%%%%%%%%%%%%%%%%%%%%%%%

\usepackage{forest}

%%%%%%%%%%%%%%%%%%%%%%%%%%%%%%%%%%
%
%	adjustbox package to force minipage alignment
%
%	Two-column structure is
%
%	\adjustbox{valign=t}
%		{
%			\begin{minipage}{0.45\textwidth}
%			...
%			\end{minipage}
%		}
%	\adjustbox{valign=t}
%		{
%		
%			\begin{minipage}{0.45\textwidth}
%			...
%			\end{minipage}
%		}
%
%%%%%%%%%%%%%%%%%%%%%%%%%%%%%%%%%%

\usepackage{adjustbox}

%%%%%%%%%%%%%%%%%%%%%%%%%%%%%%%%%
%
%		Set Beamer theme
%
%%%%%%%%%%%%%%%%%%%%%%%%%%%%%%%%%

\usetheme{Dresden}
\usecolortheme{dolphin}

%%%%%%%%%%%%%%%%%%%%%%%%%%
%
%		Hack to remove subsection bar from our
%		friends at Tex.StackExchange.Com
%
%%%%%%%%%%%%%%%%%%%%%%%%%%
\makeatletter
\beamer@theme@subsectionfalse

\makeatother
%%%%%%%%%%%%%%%%%%%%%%%%%%
%		
%		Hack to make the nice little circles appear
%		from Tex.stackexchange.com
%
%		Reference: 
%
%		https://tex.stackexchange.com/questions/2072/beamer-navigation-circles-without-subsections
%
%%%%%%%%%%%%%%%%%%%%%%%%%%

\usepackage{remreset}
\makeatletter
\@removefromreset{subsection}{section}
\makeatother
\setcounter{subsection}{1}

%%%%%%%%%%%%%%%%%%%%%%%%%%
%
%		Get rid of triangles, put in circles for 
%		itemize items
%
%%%%%%%%%%%%%%%%%%%%%%%%%%
\setbeamertemplate{items}[circle]

%%%%%%%%%%%%%%%%%%%%%%%%%%
%
%		Clean up fonts
%
%%%%%%%%%%%%%%%%%%%%%%%%%%

\usefonttheme{professionalfonts}
\usefonttheme{serif}


%%%%%%%%%%%%%%%%%%%%%%%%%%%%%%%%%
%
%		tikz plotting options
%
%%%%%%%%%%%%%%%%%%%%%%%%%%%%%%%%%

\pgfplotsset{my style/.append style={axis x line=middle, axis y line=
           middle, xlabel={}, ylabel={}, axis equal }}
           
%%%%%%%%%%%%%%%%%%%%%%%%%%%%%%%%%
%
%		Math Operators
%
%%%%%%%%%%%%%%%%%%%%%%%%%%%%%%%%%

\DeclareMathOperator{\comp}{comp}
\DeclareMathOperator{\proj}{proj}
\DeclareMathOperator{\Hess}{Hess}

\DeclareMathOperator{\curl}{curl}
\DeclareMathOperator{\diver}{div}
\DeclareMathOperator{\grad}{grad}

%%%%%%%%%%%%%%%%%%%%%%%%%%%%%%%%%
%
%	Document begins here...
%
%%%%%%%%%%%%%%%%%%%%%%%%%%%%%%%%%
           
\begin{document}

\input{213-Lecture-Macros}		%% 	macros, settings, et.c


\title{Math 213 - The Gauss Divergence Theorem}			%%	title page
\author{Peter A. Perry}
\institute{University of Kentucky}
\date{December 6, 2019}
\maketitle

%%%%%%%%%%%%%%%%%%%%%%%%%%%%%%%%%%
%	
%	Due Date Reminders
%
%%%%%%%%%%%%%%%%%%%%%%%%%%%%%%%%%%

\begin{frame}
\frametitle{Reminders}

\begin{itemize}


\item	Homework D4 (Surface integrals, Stokes' Theorem) is due tonight

\end{itemize}

\end{frame}

%%%%%%%%%%%%%%%%%%%%%%%%%%%%%%%%%%
%
%	Unit Outline
%
%%%%%%%%%%%%%%%%%%%%%%%%%%%%%%%%%%

\begin{frame}
\frametitle{Unit IV: Vector Calculus}
\begin{center}
\begin{itemize}
\item[{}]			Fundamental Theorem for Line Integrals
\item[{}]			Green's Theorem
\item[{}]			Curl and Divergence
\item[{}]			Parametric Surfaces and their Areas
\item[{}]			Surface Integrals
\item[{}]			Stokes' Theorem, I
\item[{}]			Stokes' Theorem, II 
\begingroup
\color{red}
\item[{}]			The Divergence Theorem
\endgroup
\item[{}]			
\item[{}]			Review
\item[{}]			Review
\item[{}]			Review
\end{itemize}
\end{center}
\end{frame}

\section{Learning Goals}

\begin{frame}
\frametitle{Goals of the Day}
This lecture is about the Gauss Divergence Theorem,  which illuminates the meaning of the divergence of a vector field. You will learn:

\bigskip

\begin{itemize}
\item	How the flux of a vector field over a surface bounding a \emph{simple volume} to the divergence of 
the vector field in the enclosed volume
\item	How to compute the flux of a vector field by integrating its divergence 
\end{itemize}
\end{frame}

\section{Review}

\begin{frame}
\frametitle{Vector (Differential) Calculus: The Story So Far}
\footnotesize
We have defined two `derivatives' of a vector field $\bfF$. One is a scalar and the other is a vector.

\medskip

The \emph{divergence} of a vector field 
$$\bfF = P(x,y,z) \ii + Q(x,y,z) \jj + R(x,y,z) \kk$$ 
is the scalar
$$ \mathrm{div} \, \bfF = \nabla \cdot \bfF = \frac{\dee P}{\dee x} + \frac{\dee Q}{\dee y} + \frac{\dee R}{\dee z}
$$

\vfill

The \emph{curl} of a vector field $\bfF$ is 
$$ 
\mathrm{curl} \, \bfF = \nabla \times \bfF =
\begin{vmatrix}
\ii							&		\jj							&		\kk		\\[5pt]
\dfrac{\dee}{\dee x}	&	 	\dfrac{\dee}{\dee y}	&	\dfrac{\dee}{\dee z}	\\[10pt]
P							&		Q							&	R
\end{vmatrix}
$$

\end{frame}

\begin{frame}
\frametitle{Vector (Integral) Calculus: The Story So Far}
\footnotesize
The \emph{circulation} of a vector field $\bfF$ around a closed curve $C$ is the line integral 
$$ \oint_C \bfF \cdot \, d\bfr $$
\emph{Stokes' Theorem} relates the circulation of a vector field $\bfF$ over a curve $C$ to the surface integral of its curl over any surface that bounds $C$:
$$ \iint_S \left( \nabla \times \bfF\right)  \cdot \bfn \, dS = \oint_C \bfF \cdot \, d\bfr $$

\vfill
The \emph{flux} of a vector field through a surface $S$ bounding a volume $E$is the 
surface integral
$$ \iint_S \bfF \cdot \bfn \, dS $$
where $\bfn$ is the \emph{outward} normal.  The \emph{divergence theorem}, which we'll study today, relates the flux of $\bfF$ to the integral of its divergence. 
\end{frame}

\section{The Divergence Theorem}

\begin{frame}
\frametitle{What's A \emph{Simple Volume}?}
\footnotesize
If volume $E$ is a simple volume if it has no holes and its boundary separates $\R^3$ into an ``inside'' and an ``outside.''

\medskip

\adjustbox{valign=t}
{
\begin{minipage}{0.45\textwidth}
\footnotesize
The sphere of radius $a$ centered at $(0,0,0)$

\vfill

\vfill


\begin{center}
\tdplotsetmaincoords{65}{110}
	\begin{tikzpicture}[scale=0.5,tdplot_main_coords]
	
	%%
	%%	1. 	Draw sphere
	%%
	
	\foreach \u in {7,17,...,367}
		{
			\draw[gray]
				plot[domain=0:180,variable=\v]
					(
						{3*sin(\v)*cos(\u)},
						{3*sin(\v)*sin(\u)},
						{3*cos(\v)}
					);
		}
	\foreach \v in {0,10,...,180}
		{	
			\draw[gray]
				plot[domain=0:360,variable=\u]
					(
						{3*sin(\v)*cos(\u)},
						{3*sin(\v)*sin(\u)},
						{3*cos(\v)}
					);
		}
		
		%% add red dot at $(\pi/4,\pi/4)$
		
		\draw[red,fill=red]	(
										{3*sin(45)*cos(45)},
										{3*sin(45)*sin(45)},
										{3*cos(45)}
									)
					circle(0.10cm);	
					
		%%
		%%	2. Show curve of constant v 
		%%
		
		%\uncover<2->
		{
			\draw[red,dashed]
				plot[domain=0:360,variable=\theta]
					(
						{3*sin(45)*cos(\theta)},
						{3*sin(45)*sin(\theta)},
						{3*cos(45)}
					);
		}
					
		%%
		%%	3. Show curve of constant u and redraw dot
		%%		so that it doesn't have an ugly cyan line
		%%		through it
		%%
		
		%\uncover<3->
		{
			\draw[cyan,dashed]
				plot[domain=0:180,variable=\phi]
					(
						{3*sin(\phi)*cos(45)},
						{3*sin(\phi)*sin(45)},
						{3*cos(\phi)}
					);
					
			\draw[red,fill=red]	(
										{3*sin(45)*cos(45)},
										{3*sin(45)*sin(45)},
										{3*cos(45)}
									)
					circle(0.10cm);	
		}
					
		%%
		%%	4. Add tangent vectors
		%%
		
		%\uncover<4->
		%% tangent vectors
		{
			\draw[red,very thick,->,>=stealth]
				 	(
				 			{3*sin(45)*cos(45)},
				 			{3*sin(45)*sin(45)},
				 			{3*cos(45)}
				 	)	--
				 ++(-1,1,0);
			\draw[cyan,very thick,->,>=stealth]
				 	(
				 		{3*sin(45)*cos(45)},
				 		{3*sin(45)*sin(45)},
				 		{3*cos(45)}
				 	)	--
				 ++(1,1,-{sqrt(2)});
		}
		
		
		%%
		%%	5. Add outward normal, but scale down
		%%
		
		%\uncover<6->
		{
			\draw[ultra thick,black,->,>=stealth]
					(
				 		{3*sin(45)*cos(45)},
				 		{3*sin(45)*sin(45)},
				 		{3*cos(45)}
				 	)	--	++(2,2,{2*sqrt(2)});
		}
	\end{tikzpicture}
	
	\smallskip
	
	\tiny{$\sin v \cos u \ii + \sin v \sin u \jj + \cos v \kk$}
	\end{center}
\end{minipage}
}
\quad
\adjustbox{valign=t}
{
\begin{minipage}{0.45\textwidth}
\footnotesize
The box of side $a$ in the first octant

\medskip
\tdplotsetmaincoords{70}{130}
\begin{tikzpicture}[scale=0.5,tdplot_main_coords]
	\draw[thick,->,>=stealth]	(0,0,0)	--	(5,0,0)	
		node[anchor=north east]		{\scriptsize{$x$}};
	\draw[thick,->,>=stealth]	(0,0,0)	--	(0,5,0)	
		node[anchor=west]				{\scriptsize{$y$}};
	\draw[thick,->,>=stealth]	(0,0,0)	--	(0,0,5)	
		node[anchor=south]			{\scriptsize{$z$}};
	%% xy plane face
	\draw[thick,blue]	
		(0,0,0)	--	++(3,0,0)	--	++(0,3,0)	--	++(0,-3,0)
			--	cycle;
	%%	face parallel to xz plane face but in plane
	%%	x=3
	\draw[thick,blue,fill=blue!40,opacity=0.75]
		(3,0,0)	--	++(0,3,0)	--	++(0,0,3)	--	++(0,-3,0)	
			--	cycle;
	%%  top
	\draw[thick,blue,fill=blue!20,opacity=0.75]
		(0,0,3)	--	++(3,0,0)	--	++(0,3,0)	--  ++(-3,0,0)	
			--	cycle;
	%% front side
	\draw[thick,blue,fill=blue!10,opacity=0.75]
		(3,3,0)	--	++(0,0,3)	--	++(-3,0,0)	--	++(0,0,-3)
			--	cycle;
	%% normals
	\draw[ultra thick,->,>=stealth]
		(1.5,1.5,3)	--	++(0,0,1.5);
	\draw[ultra thick,->,>=stealth]
		(1.5,3,1.5)	--	++(0,1.5,0);
	\draw[ultra thick,->,>=stealth]	
		(3,1.5,1.5)	--	++(1.5,0,0);
\end{tikzpicture}
\end{minipage}
}

\end{frame}

\begin{frame}
\frametitle{The Flux of a Vector Field out of a Box}

\footnotesize
What is the flux of a vector field
$$ \bfF(x,y,z) = P(x,y,z) \ii  + Q(x,y,z) \jj + R(x,y,z) \kk $$
out of a box of side $a$?

\medskip
There are \emph{six} surfaces, which come in pairs!

\adjustbox{valign=t}
{
	\begin{minipage}{0.45\textwidth}
	
		\scriptsize

		\medskip
		\begin{itemize}
		
		
		\bigskip

		\uncover<2->
		{
		\item[$S_1$]  $x=0$, $0 \leq y \leq a$, $0 \leq z \leq a$
		\item[$S_2$]	$x=a$, $0 \leq y \leq a$, $0 \leq z \leq a$
		}

		\bigskip
		\uncover<3->
		{
		\item[$S_3$]	$y=0$, $0 \leq x \leq a$, $0 \leq z \leq a$
		\item[$S_4$]	$y=a$,  $0 \leq x \leq a$, $0 \leq z \leq a$
		}
		
		\bigskip
		\uncover<4->
		{
		\item[$S_5$]	$z=0$, $0 \leq x \leq a$, $0 \leq y \leq a$
		\item[$S_6$]	$z=a$, $0 \leq x \leq a$, $0 \leq y \leq a$
		}

		\end{itemize}

	\end{minipage}
}
\quad
\adjustbox{valign=t}
{
	\begin{minipage}{0.45\textwidth}
	
		\tdplotsetmaincoords{60}{120}
		\begin{tikzpicture}[scale=0.6,tdplot_main_coords]
		%%
		%%	Coordinate axes
		%%
		\draw[very thick,->,>=stealth]		
				(0,0,0)	--	(5,0,0)	node[anchor=north east]	{$x$};
		\draw[very thick,->,>=stealth]		
				(0,0,0)	--	(0,5,0)	node[anchor=west]			{$y$};
		\draw[very thick,->,>=stealth]		
				(0,0,0)	--	(0,0,5)	node[anchor=south]		{$z$};
		%%
		%%	Box outline
		%%
	
		\draw[dashed,gray,opacity=0.5]
			(4,0,0)	--	
			++(0,0,4)	--	++(-4,0,0)	--	++(0,4,0)	--	
			++(0,0,-4)	--	++(4,0,0)	--	++(0,-4,0);
		\draw[dashed,gray,opacity=0.5]
			(4,0,4)	--
			++(0,4,0)	--	++	(-4,0,0);
		\draw[dashed,gray,opacity=0.5]
			(4,4,4)	--	++(0,0,-4);
		
		\only<4>
		{
		\draw[very thick,fill=blue!40,opacity=0.5]
			(0,0,0)	--	++(4,0,0)	--	++	(0,4,0)	--	++(-4,0,0)	--	++(0,-4,0);
		\draw[very thick,fill=blue!40,opacity=0.5]
			(0,0,4)	--	++(4,0,0)	--	++	(0,4,0)	--	++(-4,0,0)	--	++(0,-4,0);
		\draw[ultra thick,->,>=stealth]								(2,2,4)	--	++(0,0,1);
		\draw[ultra thick,gray,->,>=stealth,opacity=0.5]		(2,2,0)	--	++(0,0,-1);	
		\node[above] at	(2,2,0)			{$S_5$};
		\node[below]	at	(2.5,2.5,4)		{$S_6$};
		}
		
		
		\only<2>
		{
		\draw[very thick,fill=red!40,opacity=0.5]
			(0,0,0)	--	++(0,4,0)	--	++(0,0,4)	--	++(0,-4,0)	--	++(0,0,-4);
		\draw[very thick,fill=red!40,opacity=0.5]
			(4,0,0)	--	++(0,4,0)	--	++(0,0,4)	--	++(0,-4,0)	--	++(0,0,-4);
		\draw[ultra thick,->,>=stealth]
			(4,2,2)	--	++(1,0,0);
		\draw[ultra thick,gray,->,>=stealth,opacity=0.5]	
			(0,2,2)	--	++(-1,0,0);	
		\node at		(0,3,1)	{$S_1$};
		\node at		(4,3,1)	{$S_2$};
		}
		
		\only<3>
		{
		\draw[very thick,fill=purple!40,opacity=0.5]
			(0,0,0)	--	++(4,0,0)	--	++(0,0,4)	--	++(-4,0,0)	--	++(0,0,-4);
		\draw[very thick,fill=purple!40,opacity=0.5]
			(0,4,0)	--	++(4,0,0)	--	++(0,0,4)	--	++(-4,0,0)	--	++(0,0,-4);
		\draw[ultra thick,->,>=stealth]
			(2,4,2)	--	++(0,1,0);
		\draw[ultra thick,gray,->,>=stealth,opacity=0.5]
			(2,0,2)	--	++(0,-1,0);
		\node at		(2,0,1)	{$S_3$};
		\node at		(2,4,1)	{$S_4$};
		}
		
		\end{tikzpicture}
	\end{minipage}
}


\end{frame}

\begin{frame}
\frametitle{The Flux of a Vector Field out of a Box}
\footnotesize
The flux of a vector field out of a box is a sum of six terms, one for each cube face.

%\adjustbox{valign=t}
{
	\begin{minipage}{0.45\textwidth}
	
		\scriptsize
		
		\uncover<2->
		{
		\color{red}
		\begin{align*}
		&\int_0^a \int_0^a P(a,y,z) \, dy \, dz-\\
		&\int_0^a \int_0^a P(0,y,z) \, dy \, dz+
		\end{align*}
		}
		\uncover<3->
		{
		\color{purple}
		\begin{align*}
		&\int_0^a \int_0^a Q(x,a,z) \, dx \, dz-\\
		&\int_0^a \int_0^a Q(x,0,z) \, dx \, dz
		\end{align*}
		}
		\uncover<4->
		{
		\color{blue}
		\begin{align*}
		&\int_0^a \int_0^a R(x,y,a) \, dx \, dy-\\
		&\int_0^a \int_0^a R(x,y,0) \, dx \, dy+
		\end{align*}
		}
		
	\end{minipage}
}
\quad
%\adjustbox{valign=t}
{
	\begin{minipage}{0.45\textwidth}
	
		\tdplotsetmaincoords{60}{120}
		\begin{tikzpicture}[scale=0.7,tdplot_main_coords]
		%%
		%%	Coordinate axes
		%%
		\draw[very thick,->,>=stealth]		
				(0,0,0)	--	(5,0,0)	node[anchor=north east]	{$x$};
		\draw[very thick,->,>=stealth]		
				(0,0,0)	--	(0,5,0)	node[anchor=west]			{$y$};
		\draw[very thick,->,>=stealth]		
				(0,0,0)	--	(0,0,5)	node[anchor=south]		{$z$};
				
		%%
		%%	Box outline
		%%
	
		\draw[dashed,gray,opacity=0.5]
			(4,0,0)	--	
			++(0,0,4)	--	++(-4,0,0)	--	++(0,4,0)	--	
			++(0,0,-4)	--	++(4,0,0)	--	++(0,-4,0);
		\draw[dashed,gray,opacity=0.5]
			(4,0,4)	--
			++(0,4,0)	--	++	(-4,0,0);
		\draw[dashed,gray,opacity=0.5]
			(4,4,4)	--	++(0,0,-4);
		
		\only<4>
		{
		\draw[very thick,fill=blue!40,opacity=0.5]
			(0,0,0)	--	++(4,0,0)	--	++	(0,4,0)	--	++(-4,0,0)	--	++(0,-4,0);
		\draw[very thick,fill=blue!40,opacity=0.5]
			(0,0,4)	--	++(4,0,0)	--	++	(0,4,0)	--	++(-4,0,0)	--	++(0,-4,0);
		\draw[ultra thick,->,>=stealth]								(2,2,4)	--	++(0,0,1);
		\draw[ultra thick,gray,->,>=stealth,opacity=0.5]		(2,2,0)	--	++(0,0,-1);	
		\node[above] at	(2,2,0)			{$S_5$};
		\node[below]	at	(2.5,2.5,4)		{$S_6$};
		}
		
		
		\only<2>
		{
		\draw[very thick,fill=red!40,opacity=0.5]
			(0,0,0)	--	++(0,4,0)	--	++(0,0,4)	--	++(0,-4,0)	--	++(0,0,-4);
		\draw[very thick,fill=red!40,opacity=0.5]
			(4,0,0)	--	++(0,4,0)	--	++(0,0,4)	--	++(0,-4,0)	--	++(0,0,-4);
		\draw[ultra thick,->,>=stealth]
			(4,2,2)	--	++(1,0,0);
		\draw[ultra thick,gray,->,>=stealth,opacity=0.5]	
			(0,2,2)	--	++(-1,0,0);	
		\node at		(0,3,1)	{$S_1$};
		\node at		(4,3,1)	{$S_2$};
		}
		
		\only<3>
		{
		\draw[very thick,fill=purple!40,opacity=0.5]
			(0,0,0)	--	++(4,0,0)	--	++(0,0,4)	--	++(-4,0,0)	--	++(0,0,-4);
		\draw[very thick,fill=purple!40,opacity=0.5]
			(0,4,0)	--	++(4,0,0)	--	++(0,0,4)	--	++(-4,0,0)	--	++(0,0,-4);
		\draw[ultra thick,->,>=stealth]
			(2,4,2)	--	++(0,1,0);
		\draw[ultra thick,gray,->,>=stealth,opacity=0.5]
			(2,0,2)	--	++(0,-1,0);
		\node at		(2,0,1)	{$S_3$};
		\node at		(2,4,1)	{$S_4$};
		}
		
		\end{tikzpicture}
	\end{minipage}
}
\end{frame}

\begin{frame}
\frametitle{The Gauss Divergence Theorem}
\adjustbox{valign=t}
{
\begin{minipage}{0.3\textwidth}
\footnotesize
\begin{center}
\includegraphics[width=0.9\textwidth]{Images/Gauss_1828.jpg}

\smallskip
\scriptsize{Carl Friedrich Gauss, 1777-1855}
\end{center}
\end{minipage}
}
\quad
\adjustbox{valign=t}
{
\begin{minipage}{0.6\textwidth}
\footnotesize
{\color{cyan} \textbf{Theorem}} $\,$ Let $E$ be a simple solid region and let $S$ be a boundary surface of $E$, given with positive (outward) orientation. Let $\bfF$ be a vector field whose component functions have continuous partial derivatives on an open region that contains $E$. Then
Then
$$ 
\iiint_E \nabla \cdot \bfF \, dV = 
\iint_S \bfF \cdot \bfn \, dS 
$$
where $\bfn$ is the outward unit normal to $S$. 
\end{minipage}
}

\vfill

\footnotesize
Important Note: the notations 
$\bfF \cdot \bfn \, dS$ (here) and $\bfF \cdot \, d\bfS$ (the book) mean the same thing.
\end{frame}

\begin{frame}
\frametitle{The Divergence Theorem for a Cube}
\scriptsize
We can compute
$$ \iiint_V 
		\left(
			{\color{red} \frac{\dee P}{\dee x}}  	+ 
			{\color{purple} \frac{\dee Q}{\dee y}} 	+ 
			{\color{blue} \frac{\dee R}{\dee z}}
		\right)  \, dV 
$$
on a cube of side $a$ using the Fundamental Theorem of Calculus.

	\begin{minipage}{0.45\textwidth}
		\scriptsize
		\uncover<2->
		{
			\begin{align*}
			&\int_0^a \, \int_0^a \int_0^a {\color{red} \frac{\dee P}{\dee x} \, dx} \, dy \, dz\\
			&=\int_0^a \int_0^a {\color{red} \left(P(a,y,z) - P(0,y,z)\right)} \, dy \, dz 
			\end{align*}
		}
		\uncover<3->
		{
			\begin{align*}
			&\int_0^a \, \int_0^a \int_0^a {\color{purple} \frac{\dee Q}{\dee y} \,  dy} \, dx \, dz\\
			&=\int_0^a \int_0^a {\color{purple} \left(Q(x,a,z)  - Q(x,0,z)\right)} \, dx \, dz
			\end{align*}
		}
		\uncover<4->
		{
			\begin{align*}
			&\int_0^a \, \int_0^a \int_0^a {\color{blue}  \frac{\dee R}{\dee z} \, dz} \, dx \, dy\\
			&=\int_0^a \int_0^a {\color{blue} \left(R(x,y,a) - R(x,y,0)\right)} \, dx \, dy
			\end{align*}
		}
	\end{minipage}
\quad
	\begin{minipage}{0.45\textwidth}
	
		\tdplotsetmaincoords{60}{120}
		\begin{tikzpicture}[scale=0.7,tdplot_main_coords]
		%%
		%%	Coordinate axes
		%%
		\draw[very thick,->,>=stealth]		
				(0,0,0)	--	(5,0,0)	node[anchor=north east]	{$x$};
		\draw[very thick,->,>=stealth]		
				(0,0,0)	--	(0,5,0)	node[anchor=west]			{$y$};
		\draw[very thick,->,>=stealth]		
				(0,0,0)	--	(0,0,5)	node[anchor=south]		{$z$};
				
		%%
		%%	Box outline
		%%
	
		\draw[dashed,gray,opacity=0.5]
			(4,0,0)	--	
			++(0,0,4)	--	++(-4,0,0)	--	++(0,4,0)	--	
			++(0,0,-4)	--	++(4,0,0)	--	++(0,-4,0);
		\draw[dashed,gray,opacity=0.5]
			(4,0,4)	--
			++(0,4,0)	--	++	(-4,0,0);
		\draw[dashed,gray,opacity=0.5]
			(4,4,4)	--	++(0,0,-4);
		
		\only<4>
		{
		\draw[very thick,fill=blue!40,opacity=0.5]
			(0,0,0)	--	++(4,0,0)	--	++	(0,4,0)	--	++(-4,0,0)	--	++(0,-4,0);
		\draw[very thick,fill=blue!40,opacity=0.5]
			(0,0,4)	--	++(4,0,0)	--	++	(0,4,0)	--	++(-4,0,0)	--	++(0,-4,0);
		\draw[ultra thick,->,>=stealth]								(2,2,4)	--	++(0,0,1);
		\draw[ultra thick,gray,->,>=stealth,opacity=0.5]		(2,2,0)	--	++(0,0,-1);	
		\node[above] at	(2,2,0)			{$S_5$};
		\node[below]	at	(2.5,2.5,4)		{$S_6$};
		}
		
		
		\only<2>
		{
		\draw[very thick,fill=red!40,opacity=0.5]
			(0,0,0)	--	++(0,4,0)	--	++(0,0,4)	--	++(0,-4,0)	--	++(0,0,-4);
		\draw[very thick,fill=red!40,opacity=0.5]
			(4,0,0)	--	++(0,4,0)	--	++(0,0,4)	--	++(0,-4,0)	--	++(0,0,-4);
		\draw[ultra thick,->,>=stealth]
			(4,2,2)	--	++(1,0,0);
		\draw[ultra thick,gray,->,>=stealth,opacity=0.5]	
			(0,2,2)	--	++(-1,0,0);	
		\node at		(0,3,1)	{$S_1$};
		\node at		(4,3,1)	{$S_2$};
		}
		
		\only<3>
		{
		\draw[very thick,fill=purple!40,opacity=0.5]
			(0,0,0)	--	++(4,0,0)	--	++(0,0,4)	--	++(-4,0,0)	--	++(0,0,-4);
		\draw[very thick,fill=purple!40,opacity=0.5]
			(0,4,0)	--	++(4,0,0)	--	++(0,0,4)	--	++(-4,0,0)	--	++(0,0,-4);
		\draw[ultra thick,->,>=stealth]
			(2,4,2)	--	++(0,1,0);
		\draw[ultra thick,gray,->,>=stealth,opacity=0.5]
			(2,0,2)	--	++(0,-1,0);
		\node at		(2,0,1)	{$S_3$};
		\node at		(2,4,1)	{$S_4$};
		}
		
		\end{tikzpicture}
	\end{minipage}

\end{frame}

\section{Using the Divergence Theorem}
\begin{frame}
\frametitle{Using the Divergence Theorem}
\footnotesize
Compute $\int_S \bfF \cdot d \bfS$ if $$\bfF(x,y,z) = xye^z \ii + xy^2 z^3 \jj - ye^x \kk$$
and $S$ is the surface bounded by the coordinate planes and the planes $x=3$, $y=2$, and $z=1$

\medskip
\begin{minipage}{0.45\textwidth}
Using the divergence theorem we can simply integrate
${\mathrm{div}} \, \bfF$ over the region
$$\{ 0 \leq x \leq 3, \, 0 \leq y \leq 2, \, 0 \leq z \leq 1 \} $$

Set up and compute this volume integral. 

\end{minipage}
\quad
\begin{minipage}{0.45\textwidth}

\tdplotsetmaincoords{60}{120}
		\begin{tikzpicture}[scale=0.7,tdplot_main_coords]
		
		%%
		%%	Coordinate axes
		%%
		
			\draw[very thick,->,>=stealth]		
					(0,0,0)	--	(5,0,0)	node[anchor=north east]	{$x$};
			\draw[very thick,->,>=stealth]		
					(0,0,0)	--	(0,4,0)	node[anchor=west]			{$y$};
			\draw[very thick,->,>=stealth]		
					(0,0,0)	--	(0,0,3)	node[anchor=south]		{$z$};
					
		%%
		%%	The surface S
		%%
		
			\draw[very thick,fill=blue!60,opacity=0.5]
				(0,0,1)	--	++(3,0,0)	--	++(0,2,0)	--	++(-3,0,0)	--	++(0,-2,0);
			\draw[very thick,fill=blue!30,opacity=0.5]
				(3,0,0)	--	++(0,0,1)	--	++(0,2,0)	--	++(0,0,-1)	--	++(0,-2,0);
			\draw[very thick,fill=blue!30,opacity=0.5]
				(3,2,0)	--	++(0,0,1)	--	++(-3,0,0)	--	++(0,0,-1)	--	++(3,0,0);
				
		%%
		%%	Coordinates
		%%
		
			\node[below] at 	(3,0,0)						{\scriptsize{$3$}};
			\node[below] at	(0,2,0)						{\scriptsize{$2$}};
			\node[anchor=south west]	at	(0,0,1)	{\scriptsize{$1$}};
				
		\end{tikzpicture}
\end{minipage}
\end{frame}

\begin{frame}
\frametitle{Using the Divergence Theorem}
\footnotesize
\textbf{Divergence Theorem}: If $E$ is a simple closed surface and $S$ is the oriented boundary of $E$, then
$$ \iiint_E {\mathrm{div}} \,\bfF \, dV = \iint_S \bfF \cdot \, d \bfS $$

\vfill

Find $\iint_S \bfF \cdot \, d\bfS$ if 
$$ \bfF = (x^3+y^3) \ii + (y^3 + z^3) \jj + (z^3 + x^3) \kk $$
and $S$ is the sphere of radius $2$ with center at the origin.

\vfill

\begin{enumerate}
\item	Calculate ${\mathrm{div}}\, \bfF$
\medskip

\item	What's the easiest way to compute the volume integral of ${\mathrm{div}}\, \bfF$ over the sphere of radius $2$?
\end{enumerate}
\end{frame}

\begin{frame}
\frametitle{Vector Calculus Identities}
\footnotesize
\textbf{Divergence Theorem}: If $E$ is a simple closed surface and $S$ is the oriented boundary of $E$, then
$$ \iiint_E {\mathrm{div}} \,\bfF \, dV = \iint_S \bfF \cdot \, d \bfS $$

\vfill

Prove that if $\bfa$ is a constant, then
$\iint_S \bfa \cdot \, d\bfS = 0$

\vfill
Prove that $\iint_S {\mathrm{curl}}\, \bfF \cdot \, d\bfS = 0$ for a closed surface. 
(Hint: You can check that ${\mathrm{div}} \, {\mathrm{curl}}\,  \bfF = 0$).


\end{frame}

\begin{frame}
\frametitle{What We Learned About the Divergence}
\footnotesize
What does the divergence measure?  From the divergence theorem we learn
that ${\mathrm{div}}\, \bfF$ measures \emph{net outward flow per unit volume}, If 
$E$ is a very small volume surrounded by a surface $S$, then
\begin{align*}
\iiint_E {\mathrm{div}} \, \bfF \, dV 	&=	\iint_S \bfF \cdot \, d\bfS	\\[5pt]
{\mathrm{div}} \, \bfF \, \Delta V		&\simeq \iint_S \bfF \cdot d \bfS 
\end{align*}
So, for example if ${\mathrm{div}} \, \bfF = 0$, this means that the net flux is zero, i.e., $\text{inflow}  = \text{outflow}$
\end{frame}


\end{document}
